\appendix
\section{Selected Source Code Snippets}
\label{appendix:code}

This appendix provides selected code snippets from the project's source code, corresponding to the implementation details discussed in the main body of the report.

\definecolor{codegreen}{rgb}{0,0.6,0}
\definecolor{codegray}{rgb}{0.5,0.5,0.5}
\definecolor{codepurple}{rgb}{0.58,0,0.82}
\definecolor{backcolour}{rgb}{0.98,0.98,0.98}

\lstdefinestyle{mystyle}{
    backgroundcolor=\color{backcolour},   
    commentstyle=\color{codegreen},
    keywordstyle=\color{magenta},
    numberstyle=\tiny\color{codegray},
    stringstyle=\color{codepurple},
    basicstyle=\ttfamily\footnotesize,
    breakatwhitespace=false,         
    breaklines=true,                 
    captionpos=b,                    
    keepspaces=true,                 
    numbers=left,                    
    numbersep=5pt,                  
    showspaces=false,                
    showstringspaces=false,
    showtabs=false,                  
    tabsize=2
}

\lstset{style=mystyle}

\subsection{Core P2P Logic (\texttt{p2p\_daemon.py})}
\label{sec:appendix_p2p_daemon}

\subsubsection{Message Loop and Protocol Handling}
\label{subsec:appendix_message_loop}
This snippet from the \texttt{\_message\_loop} and \texttt{\_handle\_message} methods illustrates the core P2P protocol logic. It shows the idle timeout detection using \texttt{socket.timeout} and the dispatching of \texttt{PING}, \texttt{PONG}, \texttt{CLOSE}, and \texttt{CHAT} messages, as discussed in Section 3.2.1.

\begin{lstlisting}[language=Python, caption={P2P Message and Timeout Handling Logic}, label={lst:message_loop}]
def _message_loop(self, p2p_conn):
    """Read and process messages from a P2P connection."""
    # ... setup ...
    while self.running and not p2p_conn.closed:
        try:
            # Wait for data with a 1-second timeout
            chunk = p2p_conn.conn.recv(4096)
            # ... buffer processing ...
            while '\n' in buffer:
                line, buffer = buffer.split('\n', 1)
                if line.strip():
                    self._handle_message(p2p_conn, line.strip())
        
        except socket.timeout:
            # Check for idle timeout (30 seconds)
            if time.time() - p2p_conn.last_activity > self.IDLE_TIMEOUT:
                print(f"[P2P] Idle timeout for '{p2p_conn.remote_peer_id}'")
                break # Exit loop, close connection
            continue # No data, just loop again
        # ... error handling ...
    # ... cleanup ...

def _handle_message(self, p2p_conn, line):
    """Handle incoming message line."""
    try:
        msg = json.loads(line)
        msg_type = msg.get('type', 'UNKNOWN')
        
        if msg_type == 'PING':
            # Respond with PONG to keep connection alive
            self.send_message(p2p_conn.remote_peer_id, '', 'PONG')
        
        elif msg_type == 'PONG':
            # Keepalive response received, do nothing
            pass
        
        elif msg_type == 'CLOSE':
            # Graceful close requested by peer
            if self.on_message:
                self.on_message(msg.get('from'), msg.get('to'), msg)
            p2p_conn.closed = True # Mark connection for closure
        
        elif msg_type == 'CHAT':
            # Forward chat message to the application layer (HTTP bridge)
            if self.on_message:
                self.on_message(msg.get('from'), msg.get('to'), msg)
    # ... error handling ...
\end{lstlisting}

\subsubsection{Concurrency and Handshake Logic}
\label{subsec:appendix_concurrency}
The \texttt{\_accept\_loop} function (Listing \ref{lst:accept_loop}) demonstrates the concurrent server design, where each incoming connection is handled in a new thread. The \texttt{\_handle\_incoming\_connection} function (Listing \ref{lst:handshake}) shows the server-side P2P handshake protocol.

\begin{lstlisting}[language=Python, caption={Accepting and Handling New Connections}, label={lst:accept_loop}]
def _accept_loop(self):
    """Accept incoming connections."""
    while self.running:
        try:
            self.server_socket.settimeout(1.0)
            conn, addr = self.server_socket.accept()
            
            print(f"[P2P] Incoming connection from {addr}")
            
            # Handle in separate thread to not block the accept loop
            thread = threading.Thread(
                target=self._handle_incoming_connection,
                args=(conn, addr),
                daemon=True
            )
            thread.start()
            
        except socket.timeout:
            continue
        # ... error handling ...
\end{lstlisting}

\begin{lstlisting}[language=Python, caption={Server-Side Handshake Protocol}, label={lst:handshake}]
def _handle_incoming_connection(self, conn, addr):
    """Handle incoming P2P connection (server side)."""
    try:
        conn.settimeout(self.HANDSHAKE_TIMEOUT)
        
        # Read handshake line: "CONNECT <to_peer> <from_peer> <nonce>"
        handshake = self._read_line(conn)
        
        parts = handshake.strip().split()
        if len(parts) != 4 or parts[0] != 'CONNECT':
            conn.sendall(b"REJECT Invalid handshake format\n")
            conn.close()
            return
        
        _, to_peer, from_peer, nonce = parts
        
        # Verify this connection is for this peer
        if to_peer != self.peer_id:
            conn.sendall(b"REJECT Wrong peer\n")
            conn.close()
            return
        
        # Accept connection
        response = f"ACCEPT {to_peer} {from_peer} {nonce}\n"
        conn.sendall(response.encode('utf-8'))
        
        # ... create connection object and start message loop ...
\end{lstlisting}

\subsubsection{Thread Safety with RLock}
\label{subsec:appendix_thread_safety}
To ensure thread safety when accessing shared resources like the \texttt{connections} dictionary, a \texttt{threading.RLock} is used. This snippet from \texttt{send\_message} demonstrates its usage, as mentioned in Section 3.2.1.

\begin{lstlisting}[language=Python, caption={Ensuring Thread-Safe Access}, label={lst:rlock}]
class P2PDaemon:
    def __init__(self, ip, port, peer_id):
        # ...
        self.connections = {}  # remote_peer_id -> P2PConnection
        self.lock = threading.RLock() # Reentrant lock for shared state
        # ...

    def send_message(self, to_peer, body, msg_type='CHAT'):
        # The 'with' statement automatically acquires and releases the lock
        with self.lock:
            conn = self.connections.get(to_peer)
            
            if not conn:
                print(f"[P2P] No connection to '{to_peer}'")
                return False
            
            # ... build and send message ...
            
    def disconnect_peer(self, peer_id):
        with self.lock:
            if peer_id in self.connections:
                # ... send CLOSE message and remove from dictionary ...
                del self.connections[peer_id]

\end{lstlisting}

\subsection{HTTP Bridge Logic (\texttt{start\_sampleapp.py})}
\label{sec:appendix_http_bridge}

\subsubsection{HTTP to TCP Bridge (Sending Messages)}
\label{subsec:appendix_http_to_tcp}
The \texttt{/api/p2p-send} endpoint is a clear example of the "HTTP Bridge" pattern. It receives an HTTP POST request from the browser and translates it into a TCP message sent via the corresponding user's \texttt{P2PDaemon} instance.

\begin{lstlisting}[language=Python, caption={/api/p2p-send Endpoint}, label={lst:p2p_send}]
@app.route('/api/p2p-send', methods=['POST'])
def p2p_send(headers, body, username=None):
    """
    Send P2P message to a peer.
    POST /api/p2p-send with JSON: {to_peer, message, type}
    """
    try:
        data = json.loads(body) if body else {}
        to_peer = data.get('to_peer')
        message = data.get('message', '')
        
        # Get P2P daemon for the currently logged-in user
        with p2p_daemon_lock:
            daemon = p2p_daemons.get(username)
        
        if not daemon:
            return {'status': 400, 'body': '{"error": "P2P daemon not initialized."}'}
        
        # Send message via the P2P daemon's TCP socket
        success = daemon.send_message(to_peer, message, 'CHAT')
        
        if success:
            return {'status': 200, 'body': '{"status": "sent"}'}
        else:
            return {'status': 500, 'body': '{"error": "Failed to send"}'}
    # ... error handling ...
\end{lstlisting}

\subsubsection{TCP to HTTP Bridge (Receiving Messages)}
\label{subsec:appendix_tcp_to_http}
The \texttt{/api/p2p-receive} endpoint provides the bridge in the reverse direction. The \texttt{P2PDaemon} places incoming TCP messages into a user-specific queue (\texttt{message\_queues}), and this endpoint allows the browser to poll for them via HTTP.

\begin{lstlisting}[language=Python, caption={/api/p2p-receive Endpoint}, label={lst:p2p_receive}]
# Global queue for messages from TCP daemon to HTTP poller
message_queues = defaultdict(lambda: deque(maxlen=100))

# In P2PDaemon, the on_message callback does this:
def on_message(from_peer, to_peer, msg):
    with message_lock:
        message_queues[to_peer].append(msg)

# The HTTP endpoint polled by the browser:
@app.route('/api/p2p-receive', methods=['POST'])
def p2p_receive(headers, body, username=None):
    """Poll for incoming P2P messages."""
    try:
        # Get messages from this user's queue
        with message_lock:
            queue = message_queues[username]
            messages = list(queue) # Get all messages
            queue.clear() # Clear queue after reading
        
        return {
            'status': 200,
            'headers': {'Content-Type': 'application/json'},
            'body': json.dumps({'messages': messages})
        }
    # ... error handling ...
\end{lstlisting}

\subsection{State Management Logic (\texttt{tracker.py})}
\label{sec:appendix_tracker}
The \texttt{get\_expired\_peers} method is the core of the Layer 1 Heartbeat mechanism (Section 2.3.3), responsible for detecting and cleaning up peers that have not sent a heartbeat within the 45-second timeout window.

\begin{lstlisting}[language=Python, caption={Expired Peer Detection}, label={lst:expired_peers}]
class PeerTracker:
    HEARTBEAT_TIMEOUT_MS = 45000  # 45 seconds

    def get_expired_peers(self):
        """Scan for and remove expired peers."""
        now = int(time.time() * 1000)
        threshold = now - self.HEARTBEAT_TIMEOUT_MS

        expired = []
        with self.lock:
            # Iterate over a copy of items to allow modification
            for peer_id, peer in list(self.peers.items()):
                if peer.last_seen < threshold and peer.status == "ONLINE":
                    # Mark offline and notify other peers
                    peer.status = "OFFLINE"
                    expired.append(peer_id)
                    self._broadcast_event_locked('peer-left', peer)

                    # Remove from active peers map
                    del self.peers[peer_id]
        return expired
\end{lstlisting}

\subsection{Frontend Logic (\texttt{chat.js})}
\label{sec:appendix_frontend}

\subsubsection{Client-Side Polling (short-poll)}
\label{subsec:appendix_polling}
The \texttt{pollMessages} function implements the client-side polling (short-poll) mechanism. It periodically calls the \texttt{/api/p2p-receive} endpoint and processes the returned messages, including handling the special \texttt{CLOSE} message type to gracefully end a session.

\begin{lstlisting}[language=Java, caption={Polling for Messages}, label={lst:poll_messages}]
// Poll for incoming P2P messages
async function pollMessages() {
    if (!p2pRegistered) return;
    
    try {
        const response = await fetch('/api/p2p-receive', { /* ... */ });
        
        if (response.ok) {
            const data = await response.json();
            const messages = data.messages || [];
            
            messages.forEach(msg => {
                // Handle CLOSE message (session ended by remote peer)
                if (msg.type === 'CLOSE') {
                    if (msg.body === '__SESSION_ENDED__' || !msg.body) {
                        handleRemoteSessionEnd(msg.from);
                    }
                    return;
                }
                
                // Handle normal CHAT message
                if (msg.type === 'CHAT' && activePeer && msg.from === activePeer) {
                    appendMessage(msg.body, 'received', msg.from);
                }
            });
        }
    } catch (error) {
        console.error('[Chat] Error polling messages:', error);
    }
}
\end{lstlisting}

\subsubsection{Bug Fix: Handling Rejected Connections}
\label{subsec:appendix_bugfix1}
The \texttt{pollConnectionResponses} function was created to fix Bug \#1. It polls a dedicated endpoint (\texttt{/api/p2p-get-responses}) to check if a connection request has been explicitly rejected by another peer, preventing the client from waiting indefinitely.

\begin{lstlisting}[language=Java, caption={Polling for Connection Rejections}, label={lst:poll_responses}]
// Poll for connection responses (accept/reject notifications)
async function pollConnectionResponses() {
    if (!p2pRegistered) return;
    
    try {
        const response = await fetch('/api/p2p-get-responses');
        
        if (response.ok) {
            const data = await response.json();
            const responses = data.responses || [];
            
            responses.forEach(resp => {
                if (resp.status === 'rejected') {
                    // If we are waiting for this peer, show rejection message
                    if (activePeer === resp.from) {
                        appendSystemMessage(`${resp.from} rejected your connection request`, true);
                        updateConnectionStatus(false, 'Request rejected');
                        
                        // Close chat window automatically
                        setTimeout(() => {
                            // ... logic to close chat UI ...
                        }, 2000);
                    }
                }
            });
        }
    } catch (error) {
        console.error('[Chat] Error polling connection responses:', error);
    }
}
\end{lstlisting}

\section{Configuration \& Scripts}
\label{sec:appendix_config_scripts}
This section contains key configuration files and startup scripts that define how the system's components are launched and how they interact.

\subsection{Proxy Configuration (\texttt{proxy.conf})}
\label{subsec:appendix_proxy_conf}
The \texttt{proxy.conf} file configures the reverse proxy (\texttt{proxy.py}). It defines the upstream server (the main application backend) and the static file directory, routing all incoming traffic accordingly.

\begin{lstlisting}[language=bash, caption={Proxy Server Configuration}, label={lst:proxy_conf}]
# WeApRous Proxy Configuration

# Upstream application server
upstream 127.0.0.1:9001

# Directory for serving static files (HTML, CSS, JS)
static_dir ./www/
\end{lstlisting}

\subsection{Master Startup Script (\texttt{start\_all.bat})}
\label{subsec:appendix_startup_script}
The \texttt{start\_all.bat} script is the master script used to launch the entire system. It starts the backend application server and the reverse proxy in separate processes, providing a single command to get the system running.

\begin{lstlisting}[language=bash, caption={Master Startup Script for Windows}, label={lst:start_all}]
@echo off
REM
REM Copyright (C) 2025 pdnguyen of HCMC University of Technology VNU-HCM.
REM
echo "Starting WeApRous services..."

REM Start the backend server (SampleApp)
echo "Starting backend server on port 9001..."
start "Backend" cmd /c "python start_backend.py --server-port 9001"

REM Wait a moment for the backend to initialize
timeout /t 2 /nobreak > nul

REM Start the proxy server
echo "Starting proxy server on port 9000..."
start "Proxy" cmd /c "python start_proxy.py --server-port 9000"

echo "All services started."
echo "Access the application at http://127.0.0.1:9000"
\end{lstlisting}

\subsection{Application Server Entrypoint}
\label{subsec:appendix_app_entrypoint}
The main entrypoint in \texttt{start\_sampleapp.py} (aliased as \texttt{start\_backend.py}) shows how command-line arguments are parsed to configure the server's IP and port, demonstrating the application's configurability.

\begin{lstlisting}[language=Python, caption={Application Server Entrypoint Logic}, label={lst:app_entrypoint}]
if __name__ == "__main__":
    parser = argparse.ArgumentParser(
        prog='SampleApp',
        description='WeApRous sample application with authentication',
        epilog='Task 1 & 2 implementation'
    )
    parser.add_argument('--server-ip', default='127.0.0.1')
    parser.add_argument('--server-port', type=int, default=PORT)
 
    args = parser.parse_args()
    ip = args.server_ip
    port = args.server_port

    print(f"[SampleApp] Starting on {ip}:{port}")
    
    app.prepare_address(ip, port)
    app.run()
\end{lstlisting}

\section{API \& Protocol Reference}
\label{sec:appendix_api_protocol}

\subsection{HTTP API Endpoints}
\label{subsec:appendix_api}
The following table details the RESTful API endpoints provided by the HTTP Bridge (\texttt{start\_sampleapp.py}). These endpoints facilitate communication between the web browser client and the backend P2P system.

\begin{longtable}{|p{0.1\linewidth}|p{0.25\linewidth}|p{0.55\linewidth}|}
    \caption{HTTP API Endpoint Reference} \label{tab:api_reference} \\
    \hline
    \textbf{Method} & \textbf{Endpoint} & \textbf{Description} \\
    \hline
    \endfirsthead
    
    \multicolumn{3}{c}%
    {{\bfseries \tablename\ \thetable{} -- continued from previous page}} \\
    \hline
    \textbf{Method} & \textbf{Endpoint} & \textbf{Description} \\
    \hline
    \endhead

    \hline \multicolumn{3}{|r|}{{Continued on next page}} \\ \hline
    \endfoot

    \hline
    \endlastfoot

    % --- Authentication ---
    \multicolumn{3}{|c|}{\textbf{Authentication}} \\
    \hline
    POST & \texttt{/login} & Authenticates a user based on username and password. On success, sets `auth` and `username` cookies and redirects to the main page. \\
    \hline
    GET, POST & \texttt{/logout} & Clears user session cookies, unregisters the peer from the tracker, stops the P2P daemon, and redirects to the login page. \\
    \hline
    
    % --- Peer Management ---
    \multicolumn{3}{|c|}{\textbf{Peer Management \& Discovery}} \\
    \hline
    POST & \texttt{/api/submit-info} & Registers the logged-in user as a peer with the tracker. This is the first step after login to become visible to others. It also creates and starts the user's P2P daemon. \\
    \hline
    GET & \texttt{/api/get-list} & Retrieves a list of all currently online peers, excluding the requesting user. \\
    \hline
    POST & \texttt{/api/broadcast-peer} & Polling (short-poll) endpoint for receiving real-time peer events (e.g., `peer-joined`, `peer-left`). \\
    \hline
    POST & \texttt{/api/heartbeat} & Sent periodically by the client to signal to the tracker that the peer is still active. Prevents the peer from being marked as expired. \\
    \hline

    % --- Connection Lifecycle ---
    \multicolumn{3}{|c|}{\textbf{P2P Connection Lifecycle}} \\
    \hline
    POST & \texttt{/api/p2p-request} & Sends a connection request to another peer. The request is queued on the server for the target peer to retrieve. \\
    \hline
    GET & \texttt{/api/p2p-get-requests} & Polled by the client to check for any pending incoming connection requests from other peers. \\
    \hline
    POST & \texttt{/api/p2p-accept} & Accepts an incoming connection request. This triggers the server to initiate a direct TCP connection to the requesting peer. \\
    \hline
    POST & \texttt{/api/p2p-reject} & Rejects an incoming connection request. A "rejected" notification is queued for the original requester. \\
    \hline
    GET & \texttt{/api/p2p-get-responses} & Polled by the client to check for \texttt{accepted} or \texttt{rejected} responses to its own outgoing connection requests. (Fix for Bug \#1). \\
    \hline
    POST & \texttt{/api/p2p-disconnect} & Instructs the P2P daemon to close its TCP connection with a specified peer. \\
    \hline

    % --- Messaging ---
    \multicolumn{3}{|c|}{\textbf{P2P Messaging}} \\
    \hline
    POST & \texttt{/api/p2p-send} & The core of the HTTP-to-TCP bridge. Sends a chat message to a connected peer via the user's P2P daemon. \\
    \hline
    POST & \texttt{/api/p2p-receive} & The core of the TCP-to-HTTP bridge. Polls the server for any messages that the P2P daemon has received and placed in the user's message queue. \\
    \hline

    % --- Status & Debugging ---
    \multicolumn{3}{|c|}{\textbf{Status \& Debugging}} \\
    \hline
    GET & \texttt{/api/p2p-status} & Retrieves the status of the user's P2P daemon, including its running state, port, and list of active TCP connections. \\
    \hline
\end{longtable}

\subsection{P2P Messaging Protocol}
\label{subsec:appendix_p2p_protocol}
The direct peer-to-peer communication occurs over a custom, line-delimited TCP protocol.

\subsubsection{Handshake Protocol}
A handshake is required to establish a valid connection.
\begin{itemize}
    \item \textbf{Client sends:} \texttt{CONNECT <to\_peer\_id> <from\_peer\_id> <nonce>\textbackslash n}
    \begin{itemize}
        \item The client initiates a TCP connection and sends this line first.
    \end{itemize}
    \item \textbf{Server responds with ACCEPT:} \texttt{ACCEPT <to\_peer\_id> <from\_peer\_id> <nonce>\textbackslash n}
    \begin{itemize}
        \item If the \texttt{to\_peer\_id} matches the server's identity, it accepts the connection.
    \end{itemize}
    \item \textbf{Server responds with REJECT:} \texttt{REJECT <reason>\textbackslash n}
    \begin{itemize}
        \item If the handshake is malformed or the \texttt{to\_peer\_id} is incorrect, the connection is rejected and closed.
    \end{itemize}
\end{itemize}

\subsubsection{Message Frame Structure}
After a successful handshake, all subsequent communication consists of single-line JSON objects, terminated by a newline character (\texttt{\textbackslash n}). Each JSON object represents one message.

\begin{lstlisting}[caption={General P2P Message Structure}]
{
  "type": "CHAT|PING|PONG|CLOSE",
  "from": "sender_peer_id",
  "to": "receiver_peer_id",
  "timestamp": 1678886400000,
  "msg_id": "unique-uuid-for-chat-messages",
  "body": "Message content or empty"
}
\end{lstlisting}

\subsubsection{Message Types}
\begin{itemize}
    \item \textbf{CHAT}: A standard text message from one user to another. The `body` contains the message content.
    \item \textbf{PING}: A keepalive message sent periodically by the daemon to check if the connection is still alive.
    \item \textbf{PONG}: The response to a \texttt{PING} message. If a \texttt{PONG} is not received within the idle timeout (30 seconds), the connection is considered dead and is closed.
    \item \textbf{CLOSE}: A message indicating a graceful shutdown of the session. It is sent before the socket is closed, allowing the recipient's UI to be notified (Fix for Bug \#3).
\end{itemize}

\section{Detailed Test Logs}
\label{sec:test_logs}

This section presents the actual logs captured from the system during the execution of key test scenarios (as defined in Table~\ref{tab:test-scenarios}). The logs are taken directly from the backend (WeApRous Backend) and proxy (WeApRous Proxy) terminals, demonstrating that the system's behavior aligns with its intended design.

\subsection{Scenario 1: Peer Login and Registration}
\textbf{Description:} \texttt{user1} and \texttt{user2} log in. Their clients then call the \texttt{/api/submit-info} endpoint, which registers them with the tracker and starts their respective P2P daemons.

\begin{lstlisting}[language=bash, caption={Actual Logs for User Login and Registration}, label={lst:log_login_actual}]
[SampleApp] Login Success for 'user1'
[SampleApp] submit-info from user: user1
[SampleApp] P2P daemon created for 'user1' on port 9101
[Tracker] Registering peer: user1 at 127.0.0.1:9101
[Tracker] Broadcasting event: {'type': 'peer-joined', ...}
[SampleApp] Peer registered: user1 at 127.0.0.1:9101

[SampleApp] Login Success for 'user2'
[SampleApp] submit-info from user: user2
[SampleApp] P2P daemon created for 'user2' on port 9102
[Tracker] Registering peer: user2 at 127.0.0.1:9102
[Tracker] Broadcasting event: {'type': 'peer-joined', ...}
[SampleApp] Peer registered: user2 at 127.0.0.1:9102
\end{lstlisting}
\textbf{Analysis:} The logs from \texttt{WeApRous Backend} terminal confirm the sequence of events for both users. Each successful login is followed by a call to \texttt{submit-info}. The system correctly assigns unique ports (\texttt{9101} for \texttt{user1}, \texttt{9102} for \texttt{user2}) and registers them with the tracker, which in turn broadcasts \texttt{peer-joined} events. This behavior matches the logic in the \texttt{/api/submit-info} handler (see Table~\ref{tab:api_reference}).

\subsection{Scenario 2: P2P Connection Handshake (Request/Accept)}
\textbf{Description:} \texttt{user1} sends a connection request to \texttt{user2}. \texttt{user2} accepts the request, which initiates the TCP handshake.

\begin{lstlisting}[language=bash, caption={Actual Logs for P2P Handshake}, label={lst:log_handshake_actual}]
# --- Logs for user1 (Requester) ---
[SampleApp] 'user1' sent connection request to 'user2'
[P2PDaemon-user1] INFO: Accepted incoming connection from ('127.0.0.1', 52196)
[P2PDaemon-user1] DEBUG: -> RECV: HELLO user2 ...
[P2PDaemon-user1] DEBUG: <- SEND: HELLO_ACK user1
[P2PDaemon-user1] INFO: Handshake successful with user2. Connection established.
[P2P] 'user1' connected to 'user2'

# --- Logs for user2 (Acceptor) ---
[SampleApp] 'user2' polled 1 requests
[SampleApp] 'user2' accepted connection from 'user1'
[P2PDaemon-user2] INFO: Connecting to user1 at 127.0.0.1:9101...
[P2PDaemon-user2] DEBUG: <- SEND: HELLO user2 ...
[P2PDaemon-user2] DEBUG: -> RECV: HELLO_ACK user1
[P2PDaemon-user2] INFO: Handshake successful with user1. Connection established.
[P2P] 'user2' connected to 'user1'
\end{lstlisting}
\textbf{Analysis:} The logs clearly show the designed protocol flow. \texttt{user1} sends a request via the HTTP server. After \texttt{user2} accepts, \texttt{user2}'s daemon initiates the TCP connection to \texttt{user1}'s daemon. The \texttt{HELLO/HELLO\_ACK} exchange confirms a successful handshake, as implemented in the \texttt{\_handle\_incoming\_connection} method (see Listing~\ref{lst:handshake}).

\subsection{Scenario 3: P2P Chat Message Exchange}
\textbf{Description:} After connecting, \texttt{user1} sends the message "hi" to \texttt{user2}.

\begin{lstlisting}[language=bash, caption={Actual Logs for P2P Message Exchange}, label={lst:log_messaging_actual}]
# --- Logs for user1 (Sender) ---
[SampleApp] 'user1' sent CHAT to 'user2'
[P2PDaemon-user1] DEBUG: <- SEND to user2: CHAT|...|hi

# --- Logs for user2 (Receiver) ---
[P2PDaemon-user2] DEBUG: -> RECV from user1: CHAT|...|hi
[P2PDaemon-user2] DEBUG: Parsed message from user1: type=CHAT, body=hi
[P2P->Queue] Message from 'user1' queued for 'user2'
[SampleApp] 'user2' polled 1 messages
\end{lstlisting}
	extbf{Analysis:} The logs validate the HTTP-to-TCP-to-HTTP message bridge. The \texttt{/api/p2p-send} call from \texttt{user1}'s browser results in a \texttt{CHAT} message being sent over the direct TCP link. \texttt{user2}'s daemon receives it, places it in the user-specific queue, and the message is retrieved by \texttt{user2}'s browser via periodic polling to \texttt{/api/p2p-receive}. This matches the logic in Listings~\ref{lst:p2p_send} and \ref{lst:p2p_receive}.

\subsection{Scenario 4: Peer Disconnection via Heartbeat Timeout}
\textbf{Description:} \texttt{user1}'s client is closed. After 30 seconds, the tracker no longer receives heartbeats from \texttt{user1} and marks the peer as \texttt{OFFLINE}.

\begin{lstlisting}[language=bash, caption={Actual Logs for Heartbeat Timeout}, label={lst:log_timeout_actual}]
# --- Logged during user2's periodic heartbeat call ---
[SampleApp] heartbeat from user2, expired: []
...
# --- After user1 has been offline for >30 seconds ---
[SampleApp] heartbeat from user2, expired: ['user1']
[Tracker] Peer user1 expired. Last seen: ...
[Tracker] Broadcasting event: {'type': 'peer-left', 'peer': ...}
\end{lstlisting}
\textbf{Analysis:} The {WeApRous Backend} terminal log shows that \texttt{user2}'s periodic calls to \texttt{/api/heartbeat} initially report no expired peers. After the timeout period, a subsequent call correctly identifies \texttt{user1} as expired. The tracker then removes \texttt{user1} and broadcasts a \texttt{peer-left} event, demonstrating the functionality of the \texttt{get\_expired\_peers} method (see Listing~\ref{lst:expired_peers}).